\subsection{Microservicio de Reservas y Agenda}

\subsubsection{Alcance y responsabilidad}
Reservas y Agenda controla la ocupación de cada laboratorio por franja
horaria. Permite consultar la disponibilidad (RF-04), crear, consultar y
cancelar reservas propias (RF-05), detectar cruces y ofrecer alternativas
(RF-06), aplicar bloqueos efectivos en la agenda (RF-11) y registrar
asistencia a reservas confirmadas (RF-12). Laboratorios y Recursos
administra el catálogo y las restricciones operativas del laboratorio.
Identidad y Académico administra las cuentas y los permisos.

\subsubsection{Propiedad de datos}
La base de Agenda conserva las reservas y sus estados, las cancelaciones
con su motivo y usuario, los intervalos bloqueados y el registro de presencia
o ausencia de los participantes. Cada registro de asistencia y cancelación
se vincula con una reserva de la misma base. Agenda guarda identificadores de
laboratorio y usuario para referirse a datos administrados por otros servicios.

\subsubsection{Componentes internos}
La API de reservas recibe las solicitudes del Gateway. El Gestor de reservas
coordina la creación, la consulta de reservas propias y la cancelación. El
Control de agenda consulta las franjas y comprueba cruces. El Gestor de
bloqueos registra las restricciones efectivas y el Registro de asistencia
vincula cada presencia o ausencia con una reserva confirmada. Estos
componentes y sus dependencias se representan en la vista de nivel 3 del
modelo C4~\cite{brownC4Componentes}.

\subsubsection{Creación de una reserva}
Este flujo cubre la creación de una reserva
(RF-05) y la detección de conflictos de horario (RF-06).


\begin{itemize}
  \item \textbf{API de reservas.} Recibe la operación propuesta
    \texttt{POST /api/reservas} con el identificador del laboratorio,
    la fecha y las horas de inicio y fin.
  \item \textbf{Gestor de reservas.} Verifica los permisos del docente
    y consulta a Laboratorios y Recursos mediante REST síncrono para
    validar el laboratorio.
  \item \textbf{Control de agenda.} Comprueba si la franja coincide con
    una reserva o un bloqueo vigente. Si está libre, guarda la reserva
    en la base de datos propia de Agenda y devuelve su identificador.
\end{itemize}

Si la franja está ocupada, se rechaza la confirmación y se muestran
horarios alternativos cuando existan. Si Laboratorios no responde,
la solicitud queda sin confirmar.

La comprobación de cruces y la escritura se realizan en una misma
transacción, con un bloqueo por laboratorio. Las solicitudes simultáneas
para el mismo laboratorio se procesan en orden de adquisición del bloqueo.

\subsubsection{Consulta de disponibilidad}

La operación propuesta \texttt{GET /api/reservas/disponibilidad} recibe
el identificador del laboratorio, la fecha y el intervalo a consultar.
La API deriva la consulta al Control de agenda, que lee las reservas y
los bloqueos vigentes en la base de Agenda. La respuesta distingue las
franjas libres, reservadas y bloqueadas (RF-04).

La disponibilidad refleja el momento de la consulta. Al intentar crear
una reserva, el servicio vuelve a comprobar la franja dentro de la
transacción de confirmación.

\subsubsection{Consulta de reservas propias}

La operación propuesta \texttt{GET /api/reservas/mias} obtiene la identidad
del docente de la sesión autenticada. El Gestor de reservas filtra por ese
identificador en la base de Agenda y devuelve las reservas del docente con su
laboratorio, fecha, horario y estado (RF-05).

\subsubsection{Cancelación de una reserva}

La operación propuesta \texttt{POST /api/reservas/\{id\}/cancelacion} recibe
el motivo. El Gestor de reservas comprueba que la reserva pertenece al docente.
Con el mismo bloqueo por laboratorio usado en la creación, vuelve a comprobar
que sigue pendiente o confirmada. En una sola transacción cambia su estado a
cancelada y registra el motivo, la fecha y el usuario. La franja deja de estar
ocupada por esa reserva al confirmar la transacción. Los demás bloqueos y
reservas vigentes siguen contando para la disponibilidad (RF-05).

\subsubsection{Bloqueos de agenda}

La operación propuesta \texttt{POST /api/agenda/bloqueos} recibe el
laboratorio, el intervalo y el motivo. El Gestor de bloqueos valida los
permisos del administrador y registra el bloqueo en Agenda. Para coordinarse
con la creación de reservas, utiliza el mismo bloqueo por laboratorio.
Una vez confirmada la transacción, la franja deja de ofrecerse para nuevas
reservas (RF-11).

\subsubsection{Registro de asistencia}

La operación propuesta \texttt{POST /api/reservas/\{id\}/asistencia}
recibe el identificador del participante y su estado de presencia o
ausencia. El componente de asistencia comprueba el permiso del encargado
y que la reserva esté confirmada. Luego registra la asistencia en Agenda,
vinculada a la reserva y al participante (RF-12).

\subsubsection{Mecanismos de comunicación}
\paragraph{REST sobre HTTPS, síncrono.}
El Gateway remite la solicitud con el contexto de identidad autenticado.
Agenda verifica los permisos de la operación y, cuando corresponde, que la
reserva pertenezca al docente. Para validar una franja, consulta el estado
y las restricciones de Laboratorios mediante
\path{GET /laboratorios/{id}/condiciones-operativas}. La consulta es síncrona
porque la confirmación necesita conocer el resultado antes de continuar.

\paragraph{RabbitMQ sobre AMQP, asíncrono.}
El contrato propuesto por Información y Avisos contempla que Agenda
publique la confirmación, la cancelación, los conflictos detectados,
el registro de asistencia y los cambios de bloqueo efectivo. La
publicación posterior mediante un registro local de eventos permite
procesar avisos y reportes fuera de la operación de reserva. Las
notificaciones operativas de Laboratorios y los eventos de bloqueo
efectivo de Agenda se deben conciliar con Belinda y Nixon.

\subsubsection{Fallos y consistencia}
Si Laboratorios no responde, Agenda pospone la confirmación. El bloqueo
transaccional por laboratorio evita que dos escrituras concurrentes en
Agenda ocupen la misma franja. Como la consulta a Laboratorios y la
escritura en Agenda ocurren en servicios distintos, queda por acordar
el protocolo de cierre por mantenimiento y qué ocurre con las reservas
ya confirmadas.
